% !TEX root = ../../root.tex
Partial differential equations describe spatially distributed processes whose
dynamics cannot generally be captured by a finite-dimensional state without
approximation. Relevant applications include thermo-fluidic processes, such as
the digital-twin model studied by Toolhally et al.~\cite{TOOLHALLY2026103558}.
For such systems, controller synthesis must account for distributed states,
boundary conditions, and model uncertainty while retaining guarantees for the
original infinite-dimensional system.

The system-theoretic foundations for infinite-dimensional control include the
semigroup treatments of Curtain and Pritchard~\cite{curtain1978infinite} and
Curtain and Zwart~\cite{curtain1995introduction}, together with the well-posed
linear-systems framework of Staffans~\cite{staffans2005wellposed}. Early robust
control results include stabilization by finite-dimensional controllers due to
Curtain and Glover~\cite{curtain1986robust}, and the state-space treatment of
$H_\infty$ control for distributed-parameter systems by van
Keulen~\cite{vankeulen1993hinfinity}.

Several methods address analysis and control directly at the PDE level.
Krsti\'c and Smyshlyaev~\cite{krstic2008boundary} developed backstepping methods
for boundary control. Fridman and Blighovsky~\cite{fridman2012sampled} studied
robust Lyapunov-based control, while Valmorbida et
al.~\cite{valmorbida2016pdesdp} formulated semidefinite stability tests for a
class of PDEs. The PIE representation developed by
Peet~\cite{peet2021pierepresentation} and extended to input-output systems by
Shivakumar et al.~\cite{shivakumar2024GDPE} provides a bounded-operator
state-space framework for PDE analysis. Shivakumar, Das, and
Peet~\cite{shivaknew} subsequently used PIE duality to obtain convex nominal
state-feedback synthesis conditions.

Uncertainty is described here using integral quadratic constraints. The IQC
framework was introduced by Megretski and Rantzer~\cite{megretski1997system},
and its connection with storage inequalities follows the dissipativity theory
of Willems~\cite{willems1972dissipative1} and the finite-horizon formulation of
Seiler~\cite{6915700}. Within the PIE framework, Das et
al.~\cite{9303892} developed robust analysis for uncertain ODE--PDE systems,
Talitckii, Peet, and Seiler~\cite{10156335} established hard IQCs with
infinite-dimensional channels, and Lenssen et
al.~\cite{lenssen2025muanalysissynthesisframeworkpartial} developed structured
robustness, performance, and observer-synthesis conditions.

Dual IQCs have previously supported robust feedforward design in the work of
K\"ose and Scherer~\cite{kose2009robust}, and estimator and feedforward
synthesis for LPV systems in the work of Venkataraman and
Seiler~\cite{venkataraman}. Scherer~\cite{scherer2009robust} characterized
system structures for which robust synthesis is convex. For more general
dynamic multipliers, Veenman and Scherer~\cite{veenman2014iqcsynthesis} and
Wang, Pfifer, and Seiler~\cite{wang2016robustiqc} used alternating controller
and multiplier searches. Yuan and Wu~\cite{7353149} also considered
dynamic-IQC feedback for uncertain systems with time-varying state delay.

Finite-horizon use of dynamic multipliers requires particular care in the
factorization and storage terms. Carrasco and
Seiler~\cite{carrasco2018equivalence} studied doubly-hard factorizations,
Scherer and Veenman~\cite{scherer2018dynamicdissipation} introduced IQCs with
terminal costs, and Pfifer and Seiler~\cite{magicalhiddenenergy} and Hu et
al.~\cite{hu} showed how multiplier storage can reduce conservatism.

These developments provide the ingredients for robust analysis of PIE systems,
but the primal controller-synthesis condition remains nonconvex because the
controller couples with both the storage operator and the multiplier. This
paper develops an equivalent dual formulation for dynamically filtered IQCs.
The dual representation preserves the robust-performance certificate while
allowing the state-feedback controller to be synthesized through a convex LPI.
It also gives an explicit realization of the resulting controller and relates
the operator construction to finite-dimensional dynamic multipliers.
